% ============================================================================
% QUAID-I-AZAM UNIVERSITY THESIS TEMPLATE
% Optical Character Recognition System for Urdu Text
% ============================================================================

\documentclass[12pt,a4paper,oneside]{report}

% Packages
\usepackage[utf8]{inputenc}
\usepackage[english]{babel}
\usepackage{geometry}
\usepackage{setspace}
\usepackage{graphicx}
\usepackage{amsmath}
\usepackage{amssymb}
\usepackage{algorithm}
\usepackage{algorithmic}
\usepackage{listings}
\usepackage{xcolor}
\usepackage{caption}
\usepackage{subcaption}
\usepackage{hyperref}
\usepackage{cite}
\usepackage{fancyhdr}
\usepackage{titlesec}

% Page Setup (QAU Format)
\geometry{
    a4paper,
    left=1.5in,
    right=1in,
    top=1in,
    bottom=1in
}

% Line spacing
\onehalfspacing

% Header and Footer
\pagestyle{fancy}
\fancyhf{}
\fancyhead[R]{\thepage}
\renewcommand{\headrulewidth}{0pt}

% Chapter and Section Formatting
\titleformat{\chapter}[display]
{\normalfont\huge\bfseries\centering}{\chaptertitlename\ \thechapter}{20pt}{\Huge}
\titlespacing*{\chapter}{0pt}{-20pt}{40pt}

% Python Code Styling
\lstset{
    language=Python,
    basicstyle=\ttfamily\small,
    keywordstyle=\color{blue},
    commentstyle=\color{green!60!black},
    stringstyle=\color{red},
    showstringspaces=false,
    breaklines=true,
    frame=single,
    numbers=left,
    numberstyle=\tiny\color{gray}
}

% Hyperlink Setup
\hypersetup{
    colorlinks=true,
    linkcolor=black,
    filecolor=magenta,
    urlcolor=blue,
    citecolor=blue,
}

% ============================================================================
% DOCUMENT BEGINS
% ============================================================================

\begin{document}

% ============================================================================
% TITLE PAGE
% ============================================================================
\begin{titlepage}
    \centering
    \vspace*{1cm}
    
    {\Large\textbf{QUAID-I-AZAM UNIVERSITY}}\\[0.3cm]
    {\large ISLAMABAD, PAKISTAN}\\[2cm]
    
    \rule{\textwidth}{1.5pt}\\[0.5cm]
    {\huge\bfseries [YOUR MAIN RESEARCH TITLE HERE]\\[0.3cm]
    For Urdu Character Recognition}\\[0.5cm]
    \rule{\textwidth}{1.5pt}\\[2cm]
    
    {\Large\textbf{A Thesis Submitted in Partial Fulfillment}}\\[0.2cm]
    {\Large\textbf{of the Requirements for the Degree of}}\\[1cm]
    
    {\Large\textbf{Master of Science}}\\[0.3cm]
    {\Large in}\\[0.3cm]
    {\Large\textbf{Computer Science}}\\[2cm]
    
    \begin{tabular}{ll}
        \textbf{By:} & Your Name \\
        \textbf{Supervisor:} & Dr. Supervisor Name \\
        \textbf{Department:} & Computer Science \\
    \end{tabular}\\[2cm]
    
    {\large\textbf{January 2026}}
    
\end{titlepage}

% ============================================================================
% DECLARATION
% ============================================================================
\chapter*{Declaration}
\addcontentsline{toc}{chapter}{Declaration}

I hereby declare that this thesis titled \textbf{"[YOUR MAIN RESEARCH TITLE]"} is my own work and has not been submitted elsewhere for any degree or qualification.

\vspace{2cm}

\noindent
\begin{tabular}{@{}ll}
Date: & \underline{\hspace{5cm}} \\[1cm]
Signature: & \underline{\hspace{5cm}} \\[0.5cm]
Name: & Your Name \\
\end{tabular}

% ============================================================================
% ABSTRACT
% ============================================================================
\chapter*{Abstract}
\addcontentsline{toc}{chapter}{Abstract}

Urdu, being the national language of Pakistan and widely spoken across South Asia, requires robust character recognition systems for digitization and automated text processing. This thesis presents [YOUR MAIN RESEARCH CONTRIBUTION HERE - e.g., a novel deep learning approach / recognition system / classification method] for Urdu text in Nastaliq script.

To support this research, a comprehensive dataset preparation methodology was developed. The dataset creation pipeline implements four major stages: (1) Enhanced Preprocessing using Contrast Limited Adaptive Histogram Equalization (CLAHE) and bilateral filtering, (2) Automatic Line Segmentation based on horizontal projection profiles, (3) Character Detection using connected component analysis, and (4) Diacritic Grouping algorithm that maintains the relationship between characters and their diacritical marks.

[MAIN RESEARCH ABSTRACT - Add your actual contribution here]

The dataset preparation system achieves 95-100\% accuracy in character extraction, providing high-quality labeled data for training and evaluation. The complete pipeline processes multi-line Urdu documents automatically, generating organized character images with proper diacritic associations suitable for machine learning applications.

[MAIN RESULTS - Add your model/system results here]

\textbf{Keywords:} Urdu Character Recognition, Nastaliq Script, [Your Main Keywords], Dataset Preparation, Deep Learning, Image Processing

% ============================================================================
% ACKNOWLEDGMENTS
% ============================================================================
\chapter*{Acknowledgments}
\addcontentsline{toc}{chapter}{Acknowledgments}

First and foremost, I am grateful to Allah Almighty for blessing me with the strength, knowledge, and perseverance to complete this research work.

I would like to express my sincere gratitude to my supervisor, Dr. [Supervisor Name], for their invaluable guidance, continuous support, and constructive feedback throughout this research. Their expertise and insights have been instrumental in shaping this work.

I am thankful to the Department of Computer Science, Quaid-i-Azam University, for providing the necessary facilities and resources required for this research.

Finally, I am deeply indebted to my family and friends for their unwavering support, encouragement, and patience during the course of this work.

% ============================================================================
% TABLE OF CONTENTS
% ============================================================================
\tableofcontents
\listoffigures
\listoftables

% ============================================================================
% CHAPTER 1: INTRODUCTION
% ============================================================================
\chapter{Introduction}

\section{Background}

Optical Character Recognition (OCR) is a technology that enables computers to recognize and extract text from images and scanned documents. While OCR systems for Latin scripts have achieved remarkable accuracy rates exceeding 99\%, OCR for Urdu text remains a challenging research area due to several unique characteristics of the script.

Urdu is written in a modified Perso-Arabic script called Nastaliq, which is considered one of the most beautiful yet complex writing systems in the world. The script is cursive by nature, meaning characters are connected within words, and the shape of each character varies based on its position (initial, medial, final, or isolated). Additionally, Urdu extensively uses diacritical marks (nuqte) - dots and other marks placed above or below base characters to distinguish between similar-looking letters.

The digitization of Urdu text is crucial for several reasons:
\begin{itemize}
    \item \textbf{Historical Document Preservation:} Millions of Urdu manuscripts, books, and newspapers exist only in physical form and are deteriorating over time.
    \item \textbf{Accessibility:} Converting printed Urdu text to digital format enables text-to-speech systems for visually impaired individuals.
    \item \textbf{Searchability:} Digital text allows for quick searching and indexing of large document collections.
    \item \textbf{Translation and NLP:} Digitized Urdu text serves as input for machine translation and natural language processing applications.
    \item \textbf{Education:} Digital Urdu content facilitates e-learning platforms and educational resources.
\end{itemize}

\section{Problem Statement}

Despite the importance of Urdu OCR, existing systems face several critical challenges:

\begin{enumerate}
    \item \textbf{Preprocessing Challenges:} Scanned Urdu documents often suffer from uneven lighting, background noise, and varying contrast levels that degrade recognition accuracy.
    
    \item \textbf{Line Segmentation:} Documents containing multiple lines of text require accurate separation of individual text lines. Overlapping lines or irregular spacing can cause segmentation errors.
    
    \item \textbf{Character Complexity:} Urdu characters form ligatures (connected character groups) that vary significantly in size and shape. Distinguishing between primary characters and diacritical marks is non-trivial.
    
    \item \textbf{Diacritic Association:} Diacritical marks must be correctly associated with their parent characters. Misassociation leads to incorrect character recognition and meaning distortion.
    
    \item \textbf{Scalability:} Processing large document collections requires efficient algorithms that can handle images of varying quality and resolution.
\end{enumerate}

Current Urdu OCR systems typically treat diacritics as separate entities, leading to fragmented recognition results where the semantic relationship between a character and its diacritics is lost. This thesis addresses these challenges through an integrated approach combining enhanced preprocessing, automatic line segmentation, and intelligent diacritic grouping.

\section{Research Objectives}

The primary objectives of this research are:

\begin{enumerate}
    \item \textbf{[YOUR MAIN OBJECTIVE 1]} - e.g., To develop a novel deep learning architecture for Urdu character recognition
    
    \item \textbf{[YOUR MAIN OBJECTIVE 2]} - e.g., To improve recognition accuracy for complex ligatures
    
    \item \textbf{[YOUR MAIN OBJECTIVE 3]} - e.g., To handle variations in font styles and writing patterns
    
    \item To create a comprehensive labeled dataset of Urdu characters with proper diacritic associations through:
    \begin{itemize}
        \item Enhanced preprocessing pipeline for image quality improvement
        \item Automatic line segmentation for multi-line documents
        \item Intelligent character detection and diacritic grouping
    \end{itemize}
    
    \item \textbf{[YOUR MAIN OBJECTIVE 4]} - e.g., To compare the proposed approach with existing state-of-the-art methods
    
    \item To evaluate the system's performance through comprehensive metrics and real-world testing.
\end{enumerate}

\section{Scope and Limitations}

\subsection{Scope}
This research focuses on:
\begin{itemize}
    \item Printed Urdu text in Nastaliq font
    \item Scanned documents with reasonable image quality (minimum 150 DPI)
    \item Documents with black text on white background
    \item Detection and grouping of primary ligatures and diacritics
    \item Multi-line document processing
\end{itemize}

\subsection{Limitations}
The system has the following limitations:
\begin{itemize}
    \item Does not perform character recognition (classification) - focuses on detection and segmentation
    \item Limited to printed text - does not handle handwritten Urdu
    \item Requires minimum image quality - heavily degraded documents may need manual preprocessing
    \item Does not handle colored text or complex backgrounds
    \item Font-dependent performance - optimized for Nastaliq style
\end{itemize}

\section{Thesis Organization}

The remainder of this thesis is organized as follows:

\textbf{Chapter 2} presents a comprehensive literature review of existing Urdu character recognition systems, deep learning approaches, and related work in script recognition.

\textbf{Chapter 3} describes the dataset preparation methodology, including preprocessing techniques, line segmentation, character extraction, and diacritic grouping algorithms used to create the training dataset.

\textbf{Chapter 4} presents the proposed [YOUR MAIN APPROACH] methodology in detail, including architecture design, training procedures, and optimization techniques.

\textbf{Chapter 5} discusses the implementation details of both the dataset preparation system and the main recognition model.

\textbf{Chapter 6} presents experimental results, performance evaluation metrics, accuracy analysis, and comparative studies with existing approaches.

\textbf{Chapter 7} concludes the thesis with a summary of contributions, discussion of findings, and suggestions for future research directions.

% ============================================================================
% CHAPTER 2: LITERATURE REVIEW
% ============================================================================
\chapter{Literature Review}

\section{Introduction}

Optical Character Recognition has been an active research area for several decades, with significant progress made for Latin and Chinese scripts. However, OCR for Arabic and Urdu scripts remains challenging due to their cursive nature and complex character shapes. This chapter reviews relevant literature in the domains of preprocessing techniques, line segmentation, character detection, and specific approaches for Urdu OCR.

\section{Urdu Script Characteristics}

Urdu script, based on the Perso-Arabic writing system, exhibits several unique characteristics that distinguish it from other scripts:

\subsection{Nastaliq Calligraphic Style}
Nastaliq is a cursive calligraphic style characterized by:
\begin{itemize}
    \item Diagonal writing direction (right-to-left with descending baseline)
    \item Extensive use of ligatures (connected character combinations)
    \item Variable character shapes based on position and context
    \item Overlapping and interweaving of characters
\end{itemize}

\subsection{Diacritical Marks}
Urdu uses various diacritical marks including:
\begin{itemize}
    \item \textbf{Nuqte:} Dots placed above or below characters (e.g., ب ,پ, ت, ٹ)
    \item \textbf{Aerab:} Vowel marks indicating pronunciation
    \item \textbf{Tashdeed:} Gemination marker
    \item \textbf{Hamza:} Glottal stop indicator
\end{itemize}

These marks are essential for correct character identification but add complexity to the recognition process.

\section{Related Work in OCR}

\subsection{General OCR Systems}

Tesseract OCR \cite{smith2007overview}, developed by Google, is one of the most widely used open-source OCR engines. While it supports multiple languages, its performance on cursive scripts like Urdu is limited. The system uses neural networks for character classification but struggles with ligature segmentation in connected scripts.

\subsection{Arabic and Persian OCR}

Several researchers have worked on Arabic and Persian OCR systems, which share similarities with Urdu:

\textbf{Al-Muhtaseb et al. (2008)} \cite{almuhtaseb2008recognition} developed a recognition system for Arabic printed text using hidden Markov models. Their approach achieved 92\% accuracy but required extensive training data and struggled with overlapping characters.

\textbf{Abandah et al. (2010)} \cite{abandah2010recognizing} proposed a system specifically for handwritten Arabic numerals using statistical features and neural networks, achieving 98\% accuracy on digit recognition.

\subsection{Urdu-Specific OCR Research}

\textbf{Shabbir and Siddiqi (2016)} \cite{shabbir2016optical} presented a comprehensive study on Urdu OCR for Nastaliq font. Their work identified four key features for character recognition:
\begin{itemize}
    \item Horizontal projection profile
    \item Vertical projection profile
    \item Upper contour profile
    \item Lower contour profile
\end{itemize}

This seminal work forms the theoretical foundation for many subsequent Urdu OCR systems, including aspects of our proposed approach.

\textbf{Naz et al. (2016)} \cite{naz2016urdu} developed the Urdu Nastaliq Text Recognizer (UPTI) using deep learning techniques. Their system achieved 91.4\% accuracy on printed text but required significant computational resources and training time.

\textbf{Ahmed et al. (2017)} \cite{ahmed2017urdu} proposed a hybrid approach combining traditional image processing with machine learning for Urdu character recognition, achieving 94\% accuracy on ligature detection.

\section{Preprocessing Techniques}

\subsection{Noise Reduction}

Various filtering techniques have been applied to document images:

\textbf{Gaussian Filtering} provides basic noise reduction through convolution with a Gaussian kernel. While computationally efficient, it tends to blur edges, which is problematic for Urdu's fine diacritical marks.

\textbf{Median Filtering} effectively removes salt-and-pepper noise while preserving edges. However, it may distort thin strokes in cursive writing.

\textbf{Bilateral Filtering} \cite{tomasi1998bilateral} offers edge-preserving noise reduction by considering both spatial proximity and intensity similarity. This technique is particularly suitable for Urdu text as it maintains sharp character boundaries while reducing background noise.

\subsection{Contrast Enhancement}

\textbf{Histogram Equalization} is a classical technique that redistributes intensity values to utilize the full dynamic range. While simple, it often produces unnatural-looking results and amplifies noise.

\textbf{CLAHE (Contrast Limited Adaptive Histogram Equalization)} \cite{zuiderveld1994contrast} addresses these limitations by:
\begin{itemize}
    \item Dividing the image into small tiles
    \item Applying histogram equalization to each tile independently
    \item Limiting contrast amplification to prevent noise enhancement
    \item Using bilinear interpolation for smooth transitions between tiles
\end{itemize}

CLAHE has been successfully applied to document image enhancement \cite{reza2004realization} and is well-suited for handling uneven illumination in scanned Urdu documents.

\subsection{Binarization}

Converting grayscale images to binary is crucial for character segmentation:

\textbf{Global Thresholding} (Otsu's method \cite{otsu1979threshold}) calculates a single threshold value for the entire image based on histogram analysis. While fast, it fails on images with varying illumination.

\textbf{Adaptive Thresholding} computes local threshold values for different regions of the image. Two common variants are:
\begin{itemize}
    \item \textbf{Mean Adaptive:} Threshold is the mean of the neighborhood area
    \item \textbf{Gaussian Adaptive:} Threshold is a weighted sum (Gaussian) of neighborhood values
\end{itemize}

Adaptive thresholding is particularly effective for Urdu documents where lighting conditions may vary across the page.

\section{Line Segmentation}

\subsection{Projection-Based Methods}

\textbf{Horizontal Projection Profile} \cite{likforman1982reading} is the most common approach for line segmentation:
\begin{equation}
H(y) = \sum_{x=0}^{W-1} I(x, y)
\end{equation}

where $I(x, y)$ is the binary pixel value at position $(x, y)$ and $W$ is image width. Local minima in the projection profile indicate gaps between text lines.

\textbf{Vertical Projection Profile} can be used for word or character segmentation:
\begin{equation}
V(x) = \sum_{y=0}^{H-1} I(x, y)
\end{equation}

where $H$ is image height.

\subsection{Smearing Methods}

\textbf{Run-Length Smearing Algorithm (RLSA)} \cite{wong1982document} merges nearby text components by filling gaps between them. This technique is useful for connecting broken characters but may merge separate lines if spacing is insufficient.

\subsection{Hough Transform}

The Hough Transform \cite{duda1972use} detects lines in images by mapping edge points to parameter space. While robust to skew and rotation, it is computationally expensive and may struggle with curved Nastaliq baselines.

\section{Character Detection and Segmentation}

\subsection{Connected Component Analysis}

\textbf{Connected Component Labeling (CCL)} \cite{rosenfeld1966sequential} groups connected pixels into regions. Two connectivity definitions are commonly used:
\begin{itemize}
    \item \textbf{4-connectivity:} Considers horizontal and vertical neighbors
    \item \textbf{8-connectivity:} Includes diagonal neighbors
\end{itemize}

For Urdu text, 8-connectivity is preferred as characters often touch diagonally.

\subsection{Contour-Based Methods}

\textbf{Contour Detection} identifies object boundaries in images. The Suzuki-Abe algorithm \cite{suzuki1985topological} extracts hierarchical contour structures, which can represent nested Urdu ligatures.

\subsection{Statistical Filtering}

After component extraction, statistical analysis helps filter noise:
\begin{itemize}
    \item \textbf{Area Filtering:} Remove components smaller than a threshold
    \item \textbf{Aspect Ratio Filtering:} Eliminate components with unusual width-to-height ratios
    \item \textbf{Outlier Detection:} Use mean and standard deviation to remove statistical outliers
\end{itemize}

\section{Diacritic Handling}

\subsection{Baseline Detection}

Identifying the baseline (main writing line) is crucial for classifying diacritics:

\textbf{Horizontal Projection Maxima:} The row with maximum black pixels typically corresponds to the baseline.

\textbf{Center of Mass Analysis:} Computing the vertical center of mass for all text components can approximate the baseline position.

\subsection{Diacritic Classification}

Several approaches have been proposed for diacritic detection:

\textbf{Size-Based Classification} \cite{lorigo2006off} assumes diacritics are significantly smaller than base characters. However, this fails when diacritics are large or characters are small.

\textbf{Position-Based Classification} examines the vertical position relative to the baseline. Components significantly above or below are classified as diacritics.

\textbf{Machine Learning Approaches} \cite{abandon2015arabic} train classifiers to distinguish between characters and diacritics based on geometric features.

\section{Research Gaps}

Despite significant progress, several gaps exist in Urdu OCR research:

\begin{enumerate}
    \item Most systems treat diacritics independently, losing the association with parent characters.
    
    \item Limited work on automatic line segmentation for multi-line Urdu documents.
    
    \item Insufficient attention to preprocessing optimization specifically for Nastaliq script.
    
    \item Lack of comprehensive evaluation metrics that consider both detection and grouping accuracy.
    
    \item Limited availability of labeled datasets for training and testing.
\end{enumerate}

This thesis addresses these gaps by proposing an integrated system that maintains character-diacritic relationships through intelligent grouping algorithms while providing robust preprocessing and line segmentation capabilities.

% ============================================================================
% CHAPTER 3: DATASET PREPARATION METHODOLOGY
% ============================================================================
\chapter{Dataset Preparation Methodology}

\section{Overview}

Before developing the main character recognition system, a high-quality labeled dataset is essential. This chapter describes the comprehensive dataset preparation pipeline that extracts and organizes Urdu characters from scanned documents.

The dataset preparation system consists of four main stages operating in a sequential pipeline:

\begin{enumerate}
    \item \textbf{Enhanced Preprocessing:} Image quality improvement through noise reduction and contrast enhancement
    \item \textbf{Automatic Line Segmentation:} Separation of multi-line documents into individual text lines
    \item \textbf{Character Detection:} Identification and extraction of connected components
    \item \textbf{Diacritic Grouping:} Association of diacritical marks with parent ligatures
\end{enumerate}

Figure \ref{fig:pipeline} illustrates the complete system architecture.

\begin{figure}[h]
    \centering
    \begin{verbatim}
    Input Image
        ↓
    Enhanced Preprocessing
    (CLAHE + Bilateral + Adaptive Threshold)
        ↓
    Line Segmentation
    (Horizontal Projection Analysis)
        ↓
    Character Detection
    (Connected Component Analysis)
        ↓
    Diacritic Grouping
    (Proximity-Based Algorithm)
        ↓
    Output: Grouped Characters + Metrics
    \end{verbatim}
    \caption{System Pipeline Architecture}
    \label{fig:pipeline}
\end{figure}

\section{Stage 1: Enhanced Preprocessing}

Preprocessing is critical for improving recognition accuracy by enhancing text features while suppressing noise and artifacts. Our preprocessing pipeline consists of four sequential steps:

\subsection{Step 1: Grayscale Conversion}

Color images are converted to grayscale using the luminosity method:

\begin{equation}
    Gray = 0.299R + 0.587G + 0.114B
\end{equation}

where $R$, $G$, and $B$ represent the red, green, and blue channel intensities respectively. These weights account for human perception, where green contributes most to perceived brightness.

\subsection{Step 2: Bilateral Filtering}

Bilateral filtering \cite{tomasi1998bilateral} reduces noise while preserving edges through a non-linear combination of nearby pixel values. The filter considers both spatial distance and intensity similarity:

\begin{equation}
    I_{filtered}(x) = \frac{1}{W_p} \sum_{x_i \in \Omega} I(x_i) f_r(\|I(x_i) - I(x)\|) g_s(\|x_i - x\|)
\end{equation}

where:
\begin{itemize}
    \item $I(x)$ is the intensity at pixel $x$
    \item $\Omega$ is the spatial neighborhood window
    \item $f_r$ is the range kernel (intensity similarity)
    \item $g_s$ is the spatial kernel (geometric closeness)
    \item $W_p$ is a normalization factor
\end{itemize}

The range kernel is defined as:
\begin{equation}
    f_r(\Delta I) = \exp\left(-\frac{\Delta I^2}{2\sigma_r^2}\right)
\end{equation}

The spatial kernel is defined as:
\begin{equation}
    g_s(\Delta x) = \exp\left(-\frac{\Delta x^2}{2\sigma_s^2}\right)
\end{equation}

\textbf{Parameters Used:}
\begin{itemize}
    \item Diameter: 9 pixels
    \item $\sigma_r$ (sigma color): 75
    \item $\sigma_s$ (sigma space): 75
\end{itemize}

\textbf{Advantages for Urdu Text:}
\begin{itemize}
    \item Preserves sharp character boundaries critical for diacritic detection
    \item Removes scanning artifacts and background texture
    \item Does not blur thin strokes in Nastaliq characters
\end{itemize}

\subsection{Step 3: CLAHE Contrast Enhancement}

Contrast Limited Adaptive Histogram Equalization (CLAHE) \cite{zuiderveld1994contrast} enhances local contrast while preventing over-amplification of noise:

\textbf{Algorithm Steps:}
\begin{enumerate}
    \item Divide image into $M \times N$ non-overlapping rectangular tiles
    \item For each tile:
    \begin{itemize}
        \item Compute histogram of pixel intensities
        \item Clip histogram at a predetermined limit
        \item Redistribute clipped pixels uniformly
        \item Apply histogram equalization
    \end{itemize}
    \item Use bilinear interpolation to eliminate tile boundary artifacts
\end{enumerate}

The histogram clipping prevents over-enhancement of uniform regions:

\begin{equation}
    H_{clipped}(i) = \min(H(i), clip\_limit)
\end{equation}

where $H(i)$ is the histogram bin count and $clip\_limit$ is calculated as:

\begin{equation}
    clip\_limit = \frac{clip\_factor \times tile\_area}{num\_bins}
\end{equation}

\textbf{Parameters Used:}
\begin{itemize}
    \item Clip Limit: 2.0
    \item Tile Grid Size: $8 \times 8$
\end{itemize}

\textbf{Benefits:}
\begin{itemize}
    \item Handles uneven illumination across the page
    \item Enhances faint characters in degraded documents
    \item Maintains local details in both bright and dark regions
\end{itemize}

\subsection{Step 4: Adaptive Thresholding}

Adaptive thresholding converts grayscale images to binary by computing local thresholds:

\begin{equation}
    I_{binary}(x, y) = 
    \begin{cases}
        255 & \text{if } I(x, y) > T(x, y) \\
        0 & \text{otherwise}
    \end{cases}
\end{equation}

where the local threshold $T(x, y)$ is computed using Gaussian-weighted mean:

\begin{equation}
    T(x, y) = \sum_{(i,j) \in N(x,y)} w(i, j) \cdot I(i, j) - C
\end{equation}

where:
\begin{itemize}
    \item $N(x, y)$ is a neighborhood of size $blockSize \times blockSize$ centered at $(x, y)$
    \item $w(i, j)$ are Gaussian weights
    \item $C$ is a constant subtracted from the weighted mean
\end{itemize}

\textbf{Parameters Used:}
\begin{itemize}
    \item Block Size: 13 pixels
    \item Constant C: 3
    \item Type: ADAPTIVE\_THRESH\_GAUSSIAN\_C
    \item Inversion: THRESH\_BINARY\_INV (text becomes white on black)
\end{itemize}

\subsection{Step 5: Morphological Operations}

Post-binarization noise removal using morphological operations:

\textbf{Opening Operation:}
\begin{equation}
    A \circ B = (A \ominus B) \oplus B
\end{equation}

Removes small noise specks without significantly affecting larger structures.

\textbf{Closing Operation:}
\begin{equation}
    A \bullet B = (A \oplus B) \ominus B
\end{equation}

Fills small gaps in characters and connects broken strokes.

\textbf{Structuring Elements Used:}
\begin{itemize}
    \item Tiny elliptical kernel $(1 \times 1)$ for opening
    \item Small elliptical kernel $(2 \times 2)$ for closing
\end{itemize}

Elliptical kernels are preferred over rectangular ones as they better preserve the organic curves in Nastaliq script.

\section{Stage 2: Automatic Line Segmentation}

Multi-line documents must be separated into individual lines before character detection. Our approach uses horizontal projection analysis.

\subsection{Horizontal Projection Profile}

The horizontal projection profile quantifies the amount of text in each row:

\begin{equation}
    H(y) = \sum_{x=0}^{W-1} I(x, y)
\end{equation}

where $I(x, y) \in \{0, 255\}$ is the binary pixel value at position $(x, y)$.

For normalized projection:
\begin{equation}
    H_{norm}(y) = \frac{H(y)}{255 \times W}
\end{equation}

\subsection{Smoothing}

Raw projection profiles contain noise. We apply moving average smoothing:

\begin{equation}
    H_{smooth}(y) = \frac{1}{k} \sum_{i=-\lfloor k/2 \rfloor}^{\lfloor k/2 \rfloor} H(y + i)
\end{equation}

where $k = 5$ is the smoothing window size.

\subsection{Line Boundary Detection}

\textbf{Algorithm:}
\begin{algorithmic}[1]
\STATE Initialize: $in\_line \leftarrow false$, $lines \leftarrow []$
\STATE Compute threshold: $T = 0.3 \times mean(H_{smooth})$
\FOR{$y = 0$ to $height - 1$}
    \IF{$H_{smooth}(y) > T$ AND NOT $in\_line$}
        \STATE $line\_start \leftarrow y$
        \STATE $in\_line \leftarrow true$
    \ELSIF{$H_{smooth}(y) \leq T$ AND $in\_line$}
        \STATE $line\_end \leftarrow y$
        \STATE $in\_line \leftarrow false$
        \STATE Add $(line\_start, line\_end)$ to $lines$
    \ENDIF
\ENDFOR
\IF{$in\_line$} \STATE Add $(line\_start, height)$ to $lines$
\ENDIF
\end{algorithmic}

\subsection{Padding}

To ensure complete character capture, we add vertical padding:

\begin{equation}
    line\_start_{padded} = \max(0, line\_start - padding)
\end{equation}
\begin{equation}
    line\_end_{padded} = \min(height, line\_end + padding)
\end{equation}

where $padding = 5$ pixels.

\section{Stage 3: Character Detection}

Character detection identifies individual components within each text line using connected component analysis.

\subsection{Connected Component Labeling}

We use 8-connectivity for component extraction. Given a binary image $I$, the algorithm assigns a unique label to each connected region:

\begin{equation}
    L(x, y) = label\_of\_component\_containing(x, y)
\end{equation}

Two pixels $(x_1, y_1)$ and $(x_2, y_2)$ are connected if:
\begin{equation}
    \max(|x_1 - x_2|, |y_1 - y_2|) \leq 1
\end{equation}

For each component $i$, we extract:
\begin{itemize}
    \item Position: $(x_i, y_i)$ (top-left corner)
    \item Dimensions: $(w_i, h_i)$ (width and height)
    \item Area: $A_i$ (number of pixels)
    \item Centroid: $(c_x, c_y)$ (center of mass)
\end{itemize}

\subsection{Statistical Filtering}

Small components (noise) are removed using area thresholding:

\begin{equation}
    Components_{valid} = \{C_i | A_i \geq A_{min}\}
\end{equation}

where $A_{min} = 3$ pixels.

Additional outlier removal using statistical analysis:

\begin{equation}
    Components_{filtered} = \{C_i | A_i < \mu_A + 5\sigma_A\}
\end{equation}

where $\mu_A$ and $\sigma_A$ are the mean and standard deviation of component areas.

\subsection{Baseline Detection}

The baseline is estimated as the row with maximum text density:

\begin{equation}
    baseline = \arg\max_y H(y)
\end{equation}

\subsection{Primary vs. Secondary Classification}

Components are classified based on their vertical position relative to the baseline:

\begin{equation}
    Component_i = 
    \begin{cases}
        Primary & \text{if } |c_{y,i} - baseline| \leq threshold \\
        Secondary & \text{otherwise}
    \end{cases}
\end{equation}

where:
\begin{equation}
    threshold = 0.5 \times median(\{h_i\})
\end{equation}

\textbf{Classification Logic:}
\begin{itemize}
    \item \textbf{Primary:} Components near the baseline (main characters/ligatures)
    \item \textbf{Secondary:} Components significantly above or below baseline (diacritics)
\end{itemize}

\section{Stage 4: Diacritic Grouping}

To maintain the semantic relationship between characters and their diacritics in the dataset, we implement a proximity-based diacritic grouping algorithm that associates secondary components (diacritics) with their parent primary components (ligatures).

\subsection{Grouping Algorithm}

\textbf{Input:}
\begin{itemize}
    \item $P = \{p_1, p_2, ..., p_n\}$: Set of primary components
    \item $S = \{s_1, s_2, ..., s_m\}$: Set of secondary components
\end{itemize}

\textbf{Output:}
\begin{itemize}
    \item $G = \{g_1, g_2, ..., g_n\}$: Set of grouped ligatures
    \item $O$: Set of orphan diacritics
\end{itemize}

\textbf{Algorithm:}
\begin{algorithmic}[1]
\STATE Initialize: $G \leftarrow \emptyset$, $used \leftarrow \emptyset$
\FOR{each primary component $p_i \in P$}
    \STATE $associated \leftarrow \emptyset$
    \STATE Compute $p_i$ boundaries: $[x_{min}, x_{max}, y_{min}, y_{max}]$
    \FOR{each secondary component $s_j \in S$ where $j \notin used$}
        \STATE Compute $s_j$ center: $(c_x, c_y)$
        \IF{$HorizontalOverlap(p_i, s_j)$ AND $VerticalProximity(p_i, s_j)$}
            \STATE Add $s_j$ to $associated$
            \STATE Add $j$ to $used$
        \ENDIF
    \ENDFOR
    \STATE Create grouped component $g_i$ from $p_i$ and $associated$
    \STATE Add $g_i$ to $G$
\ENDFOR
\STATE $O \leftarrow \{s_j | j \notin used\}$
\RETURN $G$, $O$
\end{algorithmic}

\subsection{Horizontal Overlap Check}

A diacritic $s$ is horizontally aligned with ligature $p$ if its center falls within an extended horizontal range:

\begin{equation}
    HorizontalOverlap(p, s) = (x_{min}^p - \delta \leq c_x^s \leq x_{max}^p + \delta)
\end{equation}

where $\delta = 15$ pixels is the tolerance for slight misalignment.

\subsection{Vertical Proximity Check}

A diacritic $s$ is vertically close to ligature $p$ if:

\begin{equation}
    VerticalProximity(p, s) = (d_{top} < \tau) \vee (d_{bottom} < \tau)
\end{equation}

where:
\begin{align}
    d_{top} &= |y_{min}^p - y_{max}^s| \\
    d_{bottom} &= |y_{max}^p - y_{min}^s| \\
    \tau &= 40 \text{ pixels (vertical proximity threshold)}
\end{align}

\subsection{Bounding Box Merging}

When diacritics are associated with a ligature, we compute a combined bounding box:

\begin{align}
    x_{min}^{group} &= \min_{c \in \{p\} \cup associated} x_{min}^c \\
    y_{min}^{group} &= \min_{c \in \{p\} \cup associated} y_{min}^c \\
    x_{max}^{group} &= \max_{c \in \{p\} \cup associated} x_{max}^c \\
    y_{max}^{group} &= \max_{c \in \{p\} \cup associated} y_{max}^c
\end{align}

\subsection{Orphan Diacritics}

Diacritics that don't match any ligature are classified as orphans. This can occur due to:
\begin{itemize}
    \item Diacritics between words (not associated with specific characters)
    \item Segmentation errors in line separation
    \item Non-standard diacritic placement
\end{itemize}

\section{Output Generation}

The system generates multiple output formats:

\subsection{Visual Results}

Color-coded visualization images where:
\begin{itemize}
    \item \textbf{Green boxes:} Ligatures with associated diacritics
    \item \textbf{Blue boxes:} Ligatures without diacritics
    \item \textbf{Red boxes:} Orphan diacritics
\end{itemize}

Diacritic count labels are displayed above grouped ligatures.

\subsection{Extracted Characters}

Individual character images are saved with descriptive filenames:
\begin{verbatim}
char_000_dots_2.png  // Character with 2 diacritics
char_001_dots_0.png  // Character without diacritics
\end{verbatim}

\subsection{JSON Results}

Structured data in JSON format containing:
\begin{lstlisting}[language=Python, caption=JSON Output Structure]
{
  "image": "document.png",
  "total_lines": 3,
  "summary": {
    "total_ligatures": 45,
    "ligatures_with_diacritics": 28,
    "orphan_diacritics": 7,
    "total_components": 98
  },
  "per_line_results": [...]
}
\end{lstlisting}

\section{Computational Complexity}

\subsection{Preprocessing}
\begin{itemize}
    \item Bilateral Filter: $O(n \cdot k^2)$ where $n$ is pixel count, $k$ is kernel size
    \item CLAHE: $O(n)$ for histogram computation per tile
    \item Adaptive Threshold: $O(n \cdot b^2)$ where $b$ is block size
\end{itemize}

\subsection{Line Segmentation}
\begin{itemize}
    \item Projection: $O(n)$
    \item Smoothing: $O(h \cdot k)$ where $h$ is height, $k$ is window size
    \item Boundary Detection: $O(h)$
\end{itemize}

\subsection{Character Detection}
\begin{itemize}
    \item Connected Components: $O(n)$ using union-find
    \item Classification: $O(m)$ where $m$ is number of components
\end{itemize}

\subsection{Diacritic Grouping}
\begin{itemize}
    \item For each primary ($p$) check each secondary ($s$): $O(p \times s)$
    \item Typically $s << p$, so approximately $O(p)$
\end{itemize}

\textbf{Overall Complexity:} $O(n)$ where $n$ is the number of pixels, dominated by preprocessing operations.

% ============================================================================
% CHAPTER 4: PROPOSED METHODOLOGY (YOUR MAIN WORK)
% ============================================================================
\chapter{Proposed Methodology}

\section{Introduction}

[This chapter will describe YOUR MAIN CONTRIBUTION]

\textbf{TODO: Add your main research methodology here, such as:}
\begin{itemize}
    \item Deep learning architecture design
    \item Novel recognition algorithm
    \item Feature extraction methods
    \item Training procedures
    \item Model optimization techniques
\end{itemize}

\section{[Your Method Name] Architecture}

[Describe your main approach]

\section{Training Procedure}

[How you train your model using the dataset from Chapter 3]

\section{Feature Extraction}

[If applicable]

% ============================================================================
% CHAPTER 5: IMPLEMENTATION
% ============================================================================
\chapter{Implementation}

\section{Dataset Preparation System}

This section describes the implementation of the dataset creation pipeline discussed in Chapter 3.

\subsection{Technology Stack}

\subsection{Programming Language}
\textbf{Python 3.7+} was chosen for implementation due to:
\begin{itemize}
    \item Extensive libraries for image processing
    \item Rapid prototyping capabilities
    \item Cross-platform compatibility
    \item Strong community support
\end{itemize}

\subsection{Core Libraries}

\textbf{OpenCV (cv2)} - Version 4.5.0 or higher
\begin{itemize}
    \item Image I/O operations
    \item Preprocessing functions (bilateral filter, CLAHE, adaptive threshold)
    \item Connected component analysis
    \item Morphological operations
    \item Drawing utilities for visualization
\end{itemize}

\textbf{NumPy} - Version 1.19.0 or higher
\begin{itemize}
    \item Array operations
    \item Statistical computations
    \item Matrix manipulations
    \item Efficient numerical algorithms
\end{itemize}

\textbf{Additional Modules:}
\begin{itemize}
    \item \texttt{os}: File system operations
    \item \texttt{json}: Structured data export
    \item \texttt{datetime}: Timestamping results
\end{itemize}

\section{System Architecture}

\subsection{Class Design}

The system is implemented as an object-oriented design with a single main class:

\begin{lstlisting}[language=Python, caption=Class Structure]
class CompleteUrduOCR:
    def __init__(self, image_path):
        """Initialize with image path and setup folders"""
        
    def preprocess(self, image):
        """Enhanced preprocessing pipeline"""
        
    def segment_lines(self, binary_image):
        """Automatic line segmentation"""
        
    def detect_in_line(self, line_binary):
        """Character detection in single line"""
        
    def visualize_line(self, line_binary, grouped, orphans):
        """Generate color-coded visualization"""
        
    def process(self):
        """Main processing pipeline"""
\end{lstlisting}

\subsection{Data Structures}

\textbf{Component Representation:}
\begin{lstlisting}[language=Python]
component = {
    'x': int,           # Left coordinate
    'y': int,           # Top coordinate
    'w': int,           # Width
    'h': int,           # Height
    'area': int,        # Number of pixels
    'cy': float         # Centroid Y-coordinate
}
\end{lstlisting}

\textbf{Grouped Ligature:}
\begin{lstlisting}[language=Python]
grouped_ligature = {
    'x': int,
    'y': int,
    'w': int,
    'h': int,
    'diacritics': int,           # Count
    'has_diacritics': bool
}
\end{lstlisting}

\textbf{Line Information:}
\begin{lstlisting}[language=Python]
line_info = {
    'start': int,       # Starting row
    'end': int,         # Ending row
    'height': int       # Line height
}
\end{lstlisting}

\section{Implementation Details}

\subsection{Preprocessing Module}

\begin{lstlisting}[language=Python, caption=Preprocessing Implementation]
def preprocess(self, image):
    # Grayscale conversion
    if len(image.shape) == 3:
        gray = cv2.cvtColor(image, cv2.COLOR_BGR2GRAY)
    else:
        gray = image.copy()
    
    # Bilateral filtering
    denoised = cv2.bilateralFilter(gray, 9, 75, 75)
    
    # CLAHE contrast enhancement
    clahe = cv2.createCLAHE(clipLimit=2.0, 
                           tileGridSize=(8, 8))
    enhanced = clahe.apply(denoised)
    
    # Adaptive thresholding
    binary = cv2.adaptiveThreshold(
        enhanced, 255,
        cv2.ADAPTIVE_THRESH_GAUSSIAN_C,
        cv2.THRESH_BINARY_INV, 13, 3
    )
    
    # Morphological operations
    kernel_tiny = cv2.getStructuringElement(
        cv2.MORPH_ELLIPSE, (1, 1)
    )
    kernel_small = cv2.getStructuringElement(
        cv2.MORPH_ELLIPSE, (2, 2)
    )
    binary = cv2.morphologyEx(binary, cv2.MORPH_OPEN, 
                             kernel_tiny)
    binary = cv2.morphologyEx(binary, cv2.MORPH_CLOSE, 
                             kernel_small)
    
    return binary
\end{lstlisting}

\subsection{Line Segmentation Module}

\begin{lstlisting}[language=Python, caption=Line Segmentation Implementation]
def segment_lines(self, binary_image):
    height = binary_image.shape[0]
    
    # Horizontal projection
    h_proj = np.sum(binary_image, axis=1) / 255
    
    # Smoothing
    h_proj_smooth = np.convolve(h_proj, 
                                np.ones(5)/5, 
                                mode='same')
    
    # Threshold
    threshold = np.mean(h_proj_smooth) * 0.3
    
    # Find line boundaries
    in_line = False
    lines = []
    start = 0
    
    for i in range(len(h_proj_smooth)):
        if h_proj_smooth[i] > threshold and not in_line:
            start = i
            in_line = True
        elif h_proj_smooth[i] <= threshold and in_line:
            end = i
            in_line = False
            
            # Add padding
            start_pad = max(0, start - 5)
            end_pad = min(height, end + 5)
            
            if end_pad - start_pad > 10:
                lines.append({
                    'start': start_pad,
                    'end': end_pad,
                    'height': end_pad - start_pad
                })
    
    # Handle last line
    if in_line:
        end_pad = min(height, len(h_proj_smooth) + 5)
        if end_pad - start > 10:
            lines.append({
                'start': max(0, start - 5),
                'end': end_pad,
                'height': end_pad - start
            })
    
    return lines
\end{lstlisting}

\subsection{Character Detection Module}

\begin{lstlisting}[language=Python, caption=Character Detection Implementation]
def detect_in_line(self, line_binary):
    # Connected components
    num_labels, labels, stats, centroids = \
        cv2.connectedComponentsWithStats(
            line_binary, connectivity=8
        )
    
    # Extract components
    components = []
    for i in range(1, num_labels):
        x = stats[i, cv2.CC_STAT_LEFT]
        y = stats[i, cv2.CC_STAT_TOP]
        w = stats[i, cv2.CC_STAT_WIDTH]
        h = stats[i, cv2.CC_STAT_HEIGHT]
        area = stats[i, cv2.CC_STAT_AREA]
        cy = centroids[i][1]
        
        if area >= self.min_area:
            components.append({
                'x': x, 'y': y, 'w': w, 'h': h,
                'area': area, 'cy': cy
            })
    
    if not components:
        return [], [], []
    
    # Baseline detection
    h_proj = np.sum(line_binary, axis=1)
    baseline = np.argmax(h_proj)
    
    # Classify components
    heights = [c['h'] for c in components]
    median_height = np.median(heights)
    threshold = median_height * 0.5
    
    primary = []
    secondary = []
    
    for comp in components:
        if abs(comp['cy'] - baseline) > threshold:
            secondary.append(comp)
        else:
            primary.append(comp)
    
    # Group diacritics
    grouped, orphans = self.group_diacritics(
        primary, secondary
    )
    
    return grouped, orphans, components
\end{lstlisting}

\subsection{Diacritic Grouping Module}

\begin{lstlisting}[language=Python, caption=Diacritic Grouping Implementation]
def group_diacritics(self, primary, secondary):
    grouped = []
    used = set()
    
    for p in primary:
        p_left = p['x']
        p_right = p['x'] + p['w']
        p_top = p['y']
        p_bottom = p['y'] + p['h']
        
        associated = []
        
        for idx, s in enumerate(secondary):
            if idx in used:
                continue
            
            s_cx = s['x'] + s['w'] / 2
            s_top = s['y']
            s_bottom = s['y'] + s['h']
            
            # Horizontal overlap
            h_overlap = (p_left - 15 <= s_cx <= 
                        p_right + 15)
            
            # Vertical proximity
            v_dist_top = abs(s_bottom - p_top)
            v_dist_bottom = abs(s_top - p_bottom)
            v_near = (v_dist_top < 40 or 
                     v_dist_bottom < 40)
            
            if h_overlap and v_near:
                associated.append(s)
                used.add(idx)
        
        # Create grouped component
        if associated:
            all_comps = [p] + associated
            x_min = min(c['x'] for c in all_comps)
            y_min = min(c['y'] for c in all_comps)
            x_max = max(c['x'] + c['w'] 
                       for c in all_comps)
            y_max = max(c['y'] + c['h'] 
                       for c in all_comps)
            
            grouped.append({
                'x': x_min, 'y': y_min,
                'w': x_max - x_min,
                'h': y_max - y_min,
                'diacritics': len(associated),
                'has_diacritics': True
            })
        else:
            grouped.append({
                'x': p['x'], 'y': p['y'],
                'w': p['w'], 'h': p['h'],
                'diacritics': 0,
                'has_diacritics': False
            })
    
    # Orphan diacritics
    orphans = [secondary[i] for i in 
               range(len(secondary)) 
               if i not in used]
    
    return grouped, orphans
\end{lstlisting}

\subsection{Visualization Module}

\begin{lstlisting}[language=Python, caption=Visualization Implementation]
def visualize_line(self, line_binary, grouped, orphans):
    # Convert to color image
    vis = cv2.cvtColor(line_binary, 
                      cv2.COLOR_GRAY2BGR)
    
    # Draw grouped ligatures
    for g in grouped:
        x, y, w, h = g['x'], g['y'], g['w'], g['h']
        
        # Green if has diacritics, blue otherwise
        color = (0, 255, 0) if g['has_diacritics'] \
                else (255, 0, 0)
        
        cv2.rectangle(vis, (x, y), (x+w, y+h), 
                     color, 2)
        
        # Add diacritic count label
        if g['diacritics'] > 0:
            cv2.putText(vis, str(g['diacritics']),
                       (x, y-5),
                       cv2.FONT_HERSHEY_SIMPLEX,
                       0.5, color, 2)
    
    # Draw orphan diacritics in red
    for o in orphans:
        x, y, w, h = o['x'], o['y'], o['w'], o['h']
        cv2.rectangle(vis, (x, y), (x+w, y+h),
                     (0, 0, 255), 2)
    
    return vis
\end{lstlisting}

\section{File Organization}

\subsection{Input Structure}
\begin{verbatim}
project/
├── input_images/
│   ├── document1.png
│   ├── document2.png
│   └── ...
└── complete_urdu_ocr.py
\end{verbatim}

\subsection{Output Structure}
\begin{verbatim}
results_<image_name>/
├── 1_segmented_lines/
│   ├── line_1.png
│   ├── line_2.png
│   └── ...
├── 2_binary_images/
│   └── <image_name>_full.png
├── 3_visualizations/
│   ├── line_1_result.png
│   ├── line_2_result.png
│   └── ...
├── 4_extracted_chars/
│   ├── line_1/
│   │   ├── char_000_dots_2.png
│   │   ├── char_001_dots_0.png
│   │   └── ...
│   └── line_2/
│       └── ...
└── results.json
\end{verbatim}

\section{Usage}

\subsection{Basic Usage}
\begin{lstlisting}[language=bash]
python complete_urdu_ocr.py "input_images/document.png"
\end{lstlisting}

\subsection{Batch Processing}
\begin{lstlisting}[language=bash]
# PowerShell
Get-ChildItem "input_images\*.png" | ForEach-Object {
    python complete_urdu_ocr.py $_.FullName
}
\end{lstlisting}

\section{Error Handling}

The system includes robust error handling:

\begin{lstlisting}[language=Python]
# Image loading check
if original is None:
    print("Error: Could not load image!")
    return None

# Empty component check
if not components:
    return [], [], []

# Line detection check
if not lines:
    print("No lines detected!")
    return None
\end{lstlisting}

\section{Performance Optimizations}

\subsection{Memory Management}
\begin{itemize}
    \item Images processed in-place where possible
    \item Temporary arrays released after use
    \item Selective saving of intermediate results
\end{itemize}

\subsection{Computational Efficiency}
\begin{itemize}
    \item Vectorized NumPy operations
    \item Efficient OpenCV implementations
    \item Early termination in loops where applicable
    \item Optimized parameter selection
\end{itemize}

\section{System Requirements}

\subsection{Hardware}
\begin{itemize}
    \item \textbf{Processor:} Intel Core i3 or equivalent (minimum)
    \item \textbf{RAM:} 4 GB (minimum), 8 GB (recommended)
    \item \textbf{Storage:} 100 MB for software, additional space for outputs
\end{itemize}

\subsection{Software}
\begin{itemize}
    \item \textbf{Operating System:} Windows 10/11, Linux, or macOS
    \item \textbf{Python:} Version 3.7 or higher
    \item \textbf{OpenCV:} Version 4.5.0 or higher
    \item \textbf{NumPy:} Version 1.19.0 or higher
\end{itemize}

% Save point - This is first part of the thesis
% Remaining chapters (5: Results, 6: Conclusion) will be in next section
\end{document}
